% 06_solar_system.tex
% Ch.3: Heliosphere, 'Oumuamua, Oort Cloud, Kirkwood gaps, magnetospheres

\chapter{Solar System Impedance}
\label{ch:solar_system}


\begin{objectivebox}
\begin{itemize}
    \item Map solar system boundaries (heliopause, Oort cloud) to macroscopic acoustic impedance limits within the LC vacuum.
    \item Model planetary magnetospheres as resonant impedance cavities balancing magnetic pressure against solar wind structural strain.
    \item Derive Kirkwood gaps and orbital resonances as wave-mechanic reflection matrices ($|\Gamma| \to 1$).
\end{itemize}
\end{objectivebox}

\section{Heliospheric Impedance Profile}
\label{sec:solar_impedance}

The solar wind creates a radial impedance profile extending from
the photosphere to the interstellar medium:
\begin{resultbox}{Heliospheric Impedance Profile}
\begin{equation}
    Z_{sw}(r) = \frac{Z_0}{\sqrt{1 - (\omega_p(r)/\omega)^2}},
    \quad n_e(r) = n_0 \left(\frac{1\;\text{AU}}{r}\right)^2
    \label{eq:heliospheric_Z}
\end{equation}
\end{resultbox}
The heliopause at $\sim$120~AU acts as an impedance boundary
where $n_e$ drops sharply to interstellar values.

\section{'Oumuamua (1I/2017~U1)}
\label{sec:oumuamua}

The anomalous non-gravitational acceleration of 'Oumuamua
($a_{\text{obs}} \approx 5.0 \times 10^{-6}$~m/s$^2$ at 1~AU,
scaling as $1/r^2$) is naturally explained by solar radiation
pressure on a high area-to-mass body:
\begin{resultbox}{'Oumuamua Non-Gravitational Acceleration}
\begin{equation}
    a_{\text{rad}} = \frac{L_\odot}{4\pi r^2 c} \cdot
    \frac{A}{m} \approx 4.54 \times 10^{-6}\;\text{m/s}^2
    \quad (A/m \approx 1\;\text{m}^2/\text{kg})
    \label{eq:oumuamua_accel}
\end{equation}
\end{resultbox}
This gives a 91\% match with zero free parameters.

\section{Solar Axiom-4 Onset Radius (Internal-Field Keying)}
\label{sec:oort_cloud}
% 🔴 RETRACTED + RELABELED 2026-08-03 (Oort containment retraction -- propagation of the merged 2026-07-19 deep-space reactive-bulk MECHANISM-RETRACTED ruling into Vol 3, which never received the banner Vol 1 and Vol 4 have carried since that date). Rule 12: the section BODY below is PRESERVED verbatim; only the section TITLE is struck-and-replaced, and the original title is quoted in the header note that follows. Original title, verbatim: "Oort Cloud as Saturation Boundary". The \label{sec:oort_cloud} is deliberately UNCHANGED -- but the rationale is corrected 2026-08-03 (repair pass): this comment shipped saying it was kept "so inbound \ref's keep resolving", and a two-method sweep found ZERO inbound \ref/\autoref/\cref to sec:oort_cloud anywhere in the corpus, so no link was in fact at stake for the LABEL. It is retained on the weaker and honest ground that renaming a label buys nothing and can only break a future cross-reference. The FILENAME is a different case and the link argument DOES hold there: oort-cloud-saturation-boundary.md carries four inbound markdown links that predate this lane (vol3/claim-quality.md leaf-reference row, vol3/cosmology/ch06-solar-system/index.md, vol3/cosmology/index.md, vol3/index.md) plus a row in the generated ave-kb/.index/cites.jsonl. Lockstep with the KB twin ave-kb/vol3/cosmology/ch06-solar-system/oort-cloud-saturation-boundary.md (same-date banner), whose FILENAME is likewise retained to preserve inbound links and likewise no longer describes the section's content.

\noindent\textbf{Containment claim RETRACTED (2026-08-03; propagation of the
2026-07-19 deep-space reactive-bulk ruling into Vol~3. Rule~12: the body below
is preserved verbatim; the section title was struck and is quoted above and
here --- \emph{``Oort Cloud as Saturation Boundary''}).}  The closing sentence
of this section --- \emph{``This falls within the observed Hills Cloud range
(2,000--5,000~AU inner edge), confirming that the Oort Cloud is the solar
system's impedance boundary''} --- is \textbf{RETRACTED}.  Three reasons, none
of them new to this note:

\noindent\textbf{(a) The mechanism was retracted on 2026-07-19; Vol~3 never
received the banner.}  The ruling (Vol~1
\kbleaf{vol\_1\_foundations/chapters/04\_continuum\_electrodynamics.tex}:272; KB twin
\kbleaf{ave-kb/vol1/dynamics/ch4-continuum-electrodynamics/magnetic-saturation.md}:34)
states verbatim: \emph{``(ii) any belt/Oort position derived from that
kinetic-energy loss does not survive as a bulk-dissipation result''}, and
\emph{``The Asteroid-Belt / Oort-Cloud prediction is
\textbf{MECHANISM-RETRACTED}''} --- because a sub-yield bulk resistive stall is
\textbf{Axiom-3 forbidden} (\kbleaf{common\_equations/eq\_axiom\_3.tex}:24:
``any apparent loss must be a boundary-radiation or mode-conversion channel,
\textbf{never a bulk resistive one}'').  That stall was \textbf{the only bridge
in this corpus} from a \emph{field isocline} (the locus where $g_N = a_0$) to a
\emph{population edge} (where objects actually accumulate).  With the bridge
Ax3-forbidden, $r_{sat}$ constrains a property of the solar \emph{field} and
says nothing about the location of the comet \emph{population}, so the
containment sentence has no derivation standing behind it.  Vol~1 and Vol~4
carried the banner from 2026-07-19; this section did not --- that gap, not a
new adjudication, is what this note closes.

\noindent\textbf{(b) This chapter already says so, twenty lines further down.}
The merged consistency-class note in \S``Kirkwood Gaps as Cavity Modes''
(\texttt{06\_solar\_system.tex}:205) states: \emph{``there is \textbf{no
AVE-distinct content in the Kirkwood or Oort data}.''}  That sentence and the
containment sentence cannot both stand.

\noindent\textbf{(c) The two numbers are not comparands.}  $r_{sat}$ is an
\textbf{instantaneous force-balance isocline} --- a property of the present-day
solar field, computable from $GM_\odot$ and $a_0$ alone, referencing no object.
The Hills inner edge is an \textbf{emplacement-timescale boundary}: the radius
inside which the Galactic tide and passing-star encounters have \emph{not had
time}, over $4.6$~Gyr, to torque planetesimal perihelia clear of the planetary
region --- a quantity contingent on the whole encounter history, with a
torquing rate scaling as $a^{7/2}$ in semi-major axis.  A force-balance radius
and a $4.6$-Gyr dynamical-history radius are different \emph{kinds} of
quantity.  \textbf{Agreement between them would be a coincidence, not a
confirmation}, and would be evidence for neither.

\noindent\textbf{(c-bis) The retracted sentence is also ambiguous, and fails
under both readings.}  ``Falls within the observed Hills Cloud \emph{range}
(2,000--5,000~AU \emph{inner edge})'' admits two readings.  \textbf{(A) The
bracket \emph{is} the target range} --- then $7{,}438.9$~AU \textbf{does not
fall within it}, being $\sim$1.5$\times$ above the 5,000~AU upper end, and the
sentence is false.  \textbf{(B) The bracket locates the inner edge, and
``range'' means the cloud's extent outward from it} --- the reading implied by
\kbleaf{ave-kb/vol3/claim-quality.md}:355's ``within (not on)'' phrasing; then
``within'' is satisfied, \textbf{but the claim is vacuous}, since the satisfied
condition spans more than an order of magnitude and essentially any
$a_0$-derived radius in that decade would have ``confirmed'' it.  The ambiguity
is \textbf{surfaced, not adjudicated}: it does not need adjudicating, because
the claim fails either way, and (c) disposes of it independently of both.

\noindent\textbf{The word ``observed'' in the retracted sentence is false.}  No
object has ever been observed occupying the 2,000--5,000~AU band; the bracket
is a \textbf{simulation-inferred modelling convention} whose value spans
roughly a factor of ten across authors, and it is \textbf{uncited anywhere in
this corpus} --- \texttt{manuscript/bibliography.bib} carries \textbf{zero}
\emph{oort} / \emph{hills} / \emph{comet} entries (verified two methods,
2026-08-03).  This is deliberately \textbf{not repaired by adding a literature
citation}: sourcing the bracket would sharpen the comparand mismatch in (c),
not relieve it.

\noindent\textbf{The slot is NOT refilled.}  Per substitution-not-retraction
(A47~v11b), no replacement bracket, no replacement observable and no successor
containment claim is entered here.  The reactive band-structure re-derivation
named by the 2026-07-19 ruling remains \textbf{SPEC'd, not run}.

\noindent\textbf{What $r_{sat}$ IS --- relabel, not retraction.}  The
arithmetic is untouched: $r_{sat} = \sqrt{GM_\odot/a_0} = 7{,}438.9$~AU on
canonical constants ($a_0 = c\,H_\infty/2\pi = 1.0719\times10^{-10}$~m/s$^2$;
the printed $\approx 7{,}400$~AU is that number rounded).  What it \emph{names}
is \textbf{the radius at which the Axiom-4 kernel turns on for the solar
monopole under internal-field keying} --- equivalently, \textbf{the outer edge
of the Sun's Newtonian region}.  It is a statement about the solar field and
carries no claim about where any body sits.

\noindent\textbf{...and its very existence is gated by an unadjudicated fork
(T4 keying).}  Whether that radius exists \emph{at all} depends on the
\textbf{internal-vs-total-field keying} of the Axiom-4 kernel $S(A) =
\sqrt{1-A^2}$ --- the T4 hinge, recorded verbatim at
\kbleaf{research/2026-07-10\_collapse-target-registry.md}:271--290:
\emph{``IF the kernel $S(A)=\sqrt{1-A^2}$ keys on the total local amplitude
(internal + external strain --- the unadjudicated hinge, since an external
field is also local field), THEN an external-field pre-saturation of
$\eta_{eff}$ follows structurally''}, with the arithmetic \emph{``the Galactic
external field at the solar neighborhood is $g_{ext} \approx
2.1\times10^{-10}$ --- $\approx 2.0\,a_0$ against the AVE-derived $a_0 =
1.07\times10^{-10}$''} ($2.1/1.07 = 1.96$), \emph{``or $\approx 1.8$ against
the empirical $a_0 = 1.2\times10^{-10}$''} ($2.1/1.2 = 1.75$).
\textbf{Under total-local-field keying the Sun already sits at
$g_{ext} \approx 1.8$--$2.0\,a_0$, so no saturation transition exists anywhere
in the solar system} and $r_{sat}$ names nothing physical.  The keying is
\textbf{Grant's call --- routed 2026-08-03, not resolved here}
(flag-don't-fix).
% ★ Ratio corrected 2026-08-03 (repair pass): this sentence shipped as "g_ext ≈ 2.0--2.1 a_0", which read the MANTISSA of g_ext = 2.1e-10 m/s^2 as if it were a ratio. The registry gives both ratios explicitly at :282--:283 (1.96 against the AVE a_0, 1.75 against the empirical one), and four already-merged corpus sites print 1.8--2.0. The PHYSICS DIRECTION IS UNCHANGED: the weaker ratio, 1.75, still exceeds 1, so total-field keying still removes the transition. Only the number moved.

\noindent\textbf{Cross-volume number conflict (flagged 2026-08-03, NOT
reconciled).}  Vol~3 books $r_{sat} \approx 7{,}439$~AU; Vol~4 books
$\sim$15,000~AU for the claimed same boundary
(\kbleaf{vol\_4\_engineering/chapters/11\_experimental\_falsification.tex}:370,
\kbleaf{ave-kb/vol4/claim-quality.md}:1123,
\kbleaf{ave-kb/vol4/falsification/ch11-experimental-bench-falsification/boundary-trapping-test.md}:25),
and the illustrative driver
\kbleaf{src/scripts/vol\_3\_macroscopic/simulate\_oort\_cloud\_trap.py}:76
hand-centers a Gaussian at $10^{4.2} \approx 15{,}800$~AU.  Whether these are
\textbf{one object or two} is \textbf{routed to Grant 2026-08-03} and is
deliberately \textbf{not reconciled here}.

\medskip
\noindent\textbf{THE OBJECT DISSOLVES (2026-08-04) --- the T4 fork above is
RULED: the Axiom-4 kernel keys on the TOTAL LOCAL FIELD.}  This is execution of
a ruling, not a new adjudication: the gate written into this section on
2026-08-03 now fires.  Under total-local-field keying the Sun already sits at
$g_{ext} \approx 1.8$--$2.0\,a_0$ (both ratios re-derived at execution:
$2.1/1.0719 = 1.959$ and $2.1/1.2 = 1.750$, \textbf{both} $>1$), so
\textbf{no saturation transition exists anywhere in the solar system} and
$r_{sat}$ is \textbf{not a boundary at any radius}.  \textbf{Both Oort numbers
retire}: Vol~3's $7{,}438.9$~AU and Vol~4's $\sim$15,000~AU
\emph{``Oort Cloud transition zone''} --- so the cross-volume
\textbf{one-object-or-two} question flagged below is now \textbf{moot}, not
answered.  \textbf{What survives is strictly less than a boundary:}
$\sqrt{GM_\odot/a_0} = 7438.906$~AU remains a computable isocline of the
\emph{solar monopole taken in isolation} --- an internal-field
\textbf{counterfactual/diagnostic}, since under the ruled keying the kernel
never sees the solar field in isolation.  Rule~12: nothing above or below is
edited; the slot is \textbf{not refilled}; no claim grade moves.  \emph{(The
Vol-4-side propagation of this retirement is owed and flagged, not executed
here.)}

\medskip
\noindent\emph{Preserved body (2026-08-03, Rule 12) --- superseded by the
retraction above:}

The inner boundary of the Oort Cloud corresponds to the radius where
solar gravitational acceleration reaches the MOND threshold:
\begin{resultbox}{Oort Cloud Saturation Boundary}
\begin{equation}
    g_N(r) = \frac{GM_\odot}{r^2} = a_0
    \quad\longrightarrow\quad
    r_{\text{sat}} = \sqrt{\frac{GM_\odot}{a_0}}
    \approx 7{,}400\;\text{AU}
    \label{eq:oort_saturation}
\end{equation}
\end{resultbox}
This falls within the observed Hills Cloud range (2,000--5,000~AU inner edge),
confirming that the Oort Cloud is the solar system's \textit{impedance boundary}.

\section{Kirkwood Gaps as Cavity Modes}
\label{sec:kirkwood}

The asteroid belt gaps at mean-motion resonances with Jupiter are
modelled as cavity modes in the gravitational impedance field:
$a_{\text{gap}} = a_J \cdot (q/p)^{2/3}$.  All five major gaps
(4:1, 3:1, 5:2, 7:3, 2:1) are reproduced to $< 0.3\%$ error.

\noindent\textbf{Consistency-class note (2026-07-20; ratified-by-merge
\#741).}  The $(q/p)^{2/3}$ spacing law \emph{reduces to Kepler's third
law} applied to a mean-motion resonance:
$T_{\text{ast}}/T_J = q/p \Rightarrow (a_{\text{ast}}/a_J)^{3/2} = q/p
\Rightarrow a_{\text{ast}} = a_J (q/p)^{2/3}$.  The $2/3$ is therefore the
\textbf{Kepler exponent, not a lattice number}, and the ``cavity mode in
the gravitational impedance field'' reading is a \textbf{relabel of the
Newtonian mean-motion resonance}: it fits to $< 0.3\%$ \emph{because} it is
the Newtonian resonance condition.  Applying the same knife both ways: AVE
adds nothing observable over Newton here, and Newton likewise only fits via
the resonance condition --- the two descriptions are \textbf{peer}, and
there is \textbf{no AVE-distinct content in the Kirkwood or Oort data}.  The
only $\ell_{node}$-free AVE-distinct dimensionless numbers the band map
carries ($2/\pi$, $\sqrt{2}$, $\sqrt{10/3}$) live at the threshold /
quantum scale, some 27--35 orders of magnitude away from orbital
frequencies.  Canonical:
\kbleaf{ave-kb/vol3/cosmology/ch06-solar-system/kirkwood-gaps-cavity-modes.md}:14.

\begin{figure}[H]
    \centering
    \includegraphics[width=0.95\textwidth]{solar_system_impedance.pdf}
    \caption{Solar system as impedance profile.
    \textbf{Top left:} Heliospheric electron density and plasma
    frequency (heliopause at 120~AU marked).
    \textbf{Top right:} 'Oumuamua acceleration: AVE predicts
    $4.54 \times 10^{-6}$~m/s$^2$ vs observed $5.0 \times 10^{-6}$
    (91\% match).
    % 🔴 RELABELED 2026-08-03 (Oort containment retraction, Rule 12): original caption text, verbatim, struck-and-replaced -- "\textbf{Bottom left:} Oort Cloud as the saturation ($g = a_0$) boundary at $\sim$7,400~AU." The panel plots the same curve, because the g=a_0 isocline is a solar-FIELD radius and the "Oort Cloud" identification was the retracted population claim. Same treatment as the section title above.
    % ★ CORRECTED 2026-08-03 (repair pass): this comment shipped claiming "only its label changes". THAT WAS NOT ACCURATE -- the caption's printed VALUE also moved, from $\sim$7,400~AU to $\sim$7,439~AU (the same quantity unrounded: r_sat = 7,438.9 AU on canonical constants). Two things changed in the caption, not one: the label AND the rounding of the number. The underlying arithmetic is untouched. NOTE the consequence for the artwork: the on-figure panel text cannot be regenerated from here -- research/2026-06-07_figure-audit-ledger.md:28 heads the section "provenance gap (content checklist-clean, but no generator found)", and :29 lists "solar_system_impedance" in that section's vol_3 cluster, closing "Content current; reproducibility unverifiable." (The words "orphan generator" at :29 are attached to cross_scale_universality.pdf specifically, not to this cluster; stated exactly rather than elided.) The artwork therefore almost certainly still carries BOTH the old label and the old 7,400 rounding, so the caption/artwork mismatch disclosed in the docket fragment covers the VALUE as well as the label. Disclosed, not silently reconciled.
    \textbf{Bottom left:} solar Axiom-4 onset radius --- the $g = a_0$
    isocline at $\sim$7,439~AU under internal-field keying, i.e.\ the outer
    edge of the Newtonian region.  This is a property of the solar
    \emph{field}; it is \textbf{not} an Oort-population edge (containment
    claim retracted 2026-08-03, see \S``Solar Axiom-4 Onset Radius'').
    \textbf{As of 2026-08-04 it is not a boundary either:} the T4 keying fork
    is ruled TOTAL-field, under which no saturation transition exists anywhere
    in the solar system, so the plotted radius is an internal-field
    \emph{counterfactual isocline}, not a location.
    \textbf{Bottom right:} Kirkwood gap positions, all within 0.3\%.}
    \label{fig:solar_system}
\end{figure}

\section{Planetary Magnetospheres as Impedance Cavities}
\label{sec:planetary_magnetosphere}

Each planet's magnetic dipole field creates a magnetopause:
the boundary where magnetic pressure balances solar wind
dynamic pressure.  In impedance terms:
\begin{resultbox}{Planetary Magnetopause Standoff}
\begin{equation}
    \frac{B^2(r_{mp})}{2\mu_0} = \frac{1}{2}\rho_{sw} v_{sw}^2
    \quad\longrightarrow\quad
    r_{mp} = R_p \left(\frac{B_{eq}^2}{2\mu_0 P_{sw}}\right)^{1/6}
    \label{eq:magnetopause}
\end{equation}
\end{resultbox}
The reflection coefficient $\Gamma$ at the magnetopause is computed
using the same \texttt{reflection\_\allowbreak coefficient()} function as all
other domains.

\subsection{The Plasma Standoff vs. The Gravitational Stator}

A common pedagogical risk is conflating the physical solar plasma (which interacts with the magnetosphere) with the macroscopic orbital coupling (which interacts with the gravitational LC metric). The AVE Topo-Kinematic framework explicitly decouples these "Two Winds" emitted by the Sun:

\begin{enumerate}
    \item \textbf{The Plasma Wind:} The Solar Wind is a physical, violent stream of electrically charged plasma. As derived above, it couples electromagnetically to the planetary magnetopause. The planetary magnetosphere acts purely as an acoustic electromagnetic shock absorber standing off the kinetic plasma.
    \item \textbf{The Gravitational Metric:} Gravity, however, is structurally the background scalar tension of the LC lattice. The Sun's standing orbital acoustic cavities completely ignore the planetary magnetic envelope. Instead, orbital reactive power ($VARs$) couples robustly to the sheer volumetric density of structural topological nodes: the bulk physical mass of the planet. In a macroscopic AC Motor analogy, the dense atomic lattice of the planet acts as the ``Gravitational Stator'' perfectly locking into the Sun's LC phase-rotation. (For explicit empirical validation, see the derivation of the Earth Flyby Anomaly as an Acoustic Shear Operator at the planet's solid stator boundary $R_\oplus$ in Section~\ref{sec:flyby_anomaly}).
\end{enumerate}

\textbf{Venus and Mars as the Control Group:} This structural separation is illustrated by Venus and Mars. Neither planet possesses an active intrinsic dynamo (no global magnetosphere), and thus both suffer aggressive atmospheric stripping precisely because they lack a plasma impedance cavity. Yet, their orbits obey identically precise Topo-Kinematic Keplerian mechanics as Earth and Jupiter. They suffer zero anomalous gravitational friction because the gravitational AC stator anchors exclusively to the atomic mass lattice, operating wholly independent of the magnetic planetary boundaries.

\subsection{Standing Wave Enhancement}

At the magnetopause, the reflected solar wind superposes with the
incident field.  By transmission line theory:
\begin{resultbox}{Chapman-Ferraro Boundary Enhancement}
\begin{equation}
    B_{eff} = B_{dipole} \times (1 + |\Gamma|)
\end{equation}
\end{resultbox}
For a highly reflective boundary, $B_{eff} \to 2 B_{dipole}$, giving
$4\times$ pressure enhancement.  This is the Chapman-Ferraro result
derived from first principles of impedance matching.

\subsection{Internal Plasma Pressure}

Trapped magnetospheric particles contribute pressure via the adiabatic
invariant.  The trapped fraction is derived from the dipole loss cone:
\begin{resultbox}{Dipole Loss Cone Trapped Fraction}
\begin{equation}
    f_{trap} = 1 - \sqrt{1 - (R_p/R_{mp})^3}
\end{equation}
\end{resultbox}

\subsection{Jupiter: Io Plasma Torus}

Jupiter has a dominant internal plasma source: Io's volcanic
outgassing.  The derivation chain:
\begin{align}
    Q_{tidal} &= \frac{21}{2}\frac{k_2}{Q}\left(\frac{R_{Io}}{a_{Io}}\right)^5
    \frac{n_{orb} \cdot G \cdot M_J^2 \cdot e^2}{a_{Io}} \\
    \dot{m} &= Q_{tidal} \cdot \eta_{volc} / L_{SO_2} \\
    n_{torus} &= \dot{m} / (m_{ion} \cdot V_{torus} \cdot \nu_{loss}) \\
    P_{Io} &= \frac{1}{2} n_{mp} \cdot m_{ion} \cdot (\Omega_J \cdot r_{mp})^2
\end{align}

\subsection{The Uranus Anomaly}

Uranus is unique: its magnetic dipole is tilted $59^\circ$ from
the rotation axis and offset by 0.31~$R_U$ from center.  This
creates a highly asymmetric, time-varying impedance cavity whose
magnetopause standoff varies from 14.9 to 20.8~$R_U$ as the
planet rotates (asymmetry ratio 1.40$\times$).

\subsection{Results}

\begin{center}
\small
\begin{tabularx}{\textwidth}{@{}lrrX@{}}
\toprule
\textbf{Planet} & \textbf{Predicted [R$_p$]} & \textbf{Observed [R$_p$]} & \textbf{Error} \\
\midrule
Earth   &  9.1 & 10.0 &  8.7\% \\
Jupiter & 55.6 & 63.0 & 11.8\% \\
Saturn  & 17.0 & 22.0 & 22.8\% \\
Uranus  & 22.1 & 25.0 & 11.6\% \\
Neptune & 21.7 & 26.0 & 16.4\% \\
\bottomrule
\end{tabularx}
\end{center}

\begin{figure}[H]
    \centering
    \includegraphics[width=0.95\textwidth]{planetary_magnetospheres.pdf}
    \caption{Planetary magnetospheres as impedance cavities.
    \textbf{Top left:} Uranus magnetopause cross-section showing
    the 59$^\circ$ tilted dipole, 0.31~$R_U$ offset, solar wind
    dynamic pressure $P_{sw}$ (orange arrows), and outward
    magnetic pressure $B^2/2\mu_0$ (green arrows).
    \textbf{Top right:} Reflection coefficient $\Gamma$ vs
    magnetic longitude --- the asymmetric profile oscillates
    with rotation.
    \textbf{Bottom left:} Magnetopause standoff comparison for
    all five planets (predicted vs observed).
    \textbf{Bottom right:} Dipole anomaly map --- Uranus and
    Neptune occupy the anomalous zone ($> 45^\circ$ tilt,
    $> 0.25\,R_p$ offset).}
    \label{fig:magnetospheres}
\end{figure}

\section{The Orbital LC Friction Paradox Resolved}
\label{sec:orbital_friction}

A common historical critique of continuous dielectric vacuum models (such as the rigid K4 lattice) is the Orbital Friction Paradox: if planets physically move through an inductive, reactive medium, why do their orbits not decay rapidly via inductive drag?

The physics engine resolves this purely algebraically via the Lense-Thirring strain amplitude operators (\texttt{solar\_system\_inductive\_drag.py}). By evaluating the 3D kinematic tensor $J_{spin}$ for every body in the solar system, the local gravitomagnetic phase-strain $A_{gm}$ for all celestial bodies scales around $\sim 10^{-34}$.

This establishes that the entire solar system operates deep inside \textbf{Regime~I ($A_{gm} \ll 0.121$)}. Because the strain operates infinitesimally below the Compton saturation limit, the lattice behaves entirely linearly. The frame-dragging field perfectly co-rotates spherically symmetrically, resulting in a strictly conservative, lossless induction. The phase-drag power dissipates $0$~Watts of macroscopic orbital kinetic energy. 

\begin{figure}[H]
    \centering
    \begin{minipage}{0.48\textwidth}
        \centering
        \includegraphics[width=\textwidth]{gw_power_bar.png}
    \end{minipage}\hfill
    \begin{minipage}{0.48\textwidth}
        \centering
        \includegraphics[width=\textwidth]{impedance_map_polar.png}
    \end{minipage}
    \caption{\textbf{Solar System Rotational Topology.} \textbf{Left:} Gravitational Wave baseline power emitted (conservative quadrupole radiation, distinct from the dissipative inductive drag shown identically zero at right) plotted logarithmically for all planetary bodies, normalized to Earth. \textbf{Right:} Log-scale impedance spatial mapping verifying zero non-conservative phase drag across the planetary radii.}
    \label{fig:solar_system_induction}
\end{figure}

\begin{summarybox}
\begin{itemize}
    \item Solar system boundaries, including the heliopause and the Oort cloud, map precisely to macroscopic acoustic impedance limits within the LC vacuum network.
    \item Planetary magnetospheres structurally act as resonant impedance cavities balancing local magnetic pressure against the solar wind structural strain.
    \item Orbital resonances and the Kirkwood gaps are derived directly as wave-mechanic reflection matrices ($|\Gamma| \to 1$) preventing stable accretion modes.
\end{itemize}
\end{summarybox}

\begin{exercisebox}
\begin{enumerate}
    \item Calculate the dynamic reflection coefficient $\Gamma$ for a planetary magnetopause given specifically the modeled solar wind impedance mismatch.
    \item Formulate the solar heliopause boundary strictly as an absolute LC network impedance limit for the Sun's propagating acoustic cavity.
\end{enumerate}
\end{exercisebox}
